\documentclass[a4paper,10pt]{report}\input{../0.Préambule/Préambule_cours_prof}\begin{document}

\newpage
\section*{Expressions littérales 1}
\addcontentsline{toc}{section}{Expressions littérales 1}

%\subsection*{Utiliser les conventions d'écriture}
%\addcontentsline{toc}{subsection}{Utiliser les conventions d'écriture}

\vspace{-4mm}
\begin{exer1}
Place tous les signes \og $\times$ \fg{} sous-entendus dans les expressions littérales suivantes

\vspace{-2mm}
\begin{enumerate}[font=\color{exer1}]
\begin{tabularx}{\linewidth}{*{5}{>{\arraybackslash}X}}
\item $3a$ &
\item $-5b$ &
\item $23 + 8h$ &
\item $m^2 - 5g$ &
\item $12k(g + h)$ \\
\end{tabularx}
\end{enumerate}
\end{exer1}

\vspace{-1mm}
\begin{exer1}
Simplifie les écritures littérales suivantes

\vspace{-2mm}
\begin{enumerate}[font=\color{exer1}]
\begin{tabularx}{\linewidth}{*{8}{>{\arraybackslash}X}}
\item $6 \times a$ &
\item $8 \times f$ &
\item $23 \times d$ &
\item $a \times b$ &
\item $x \times 9$ &
\item $y \times 3$ &
\item $c \times 5$ &
\item $g \times 12$ \\
\end{tabularx}
\end{enumerate}
\end{exer1}

\vspace{-1mm}
\begin{exer1}
Simplifie les écritures littérales suivantes

\vspace{-2mm}
\begin{enumerate}[font=\color{exer1}]
\begin{tabularx}{\linewidth}{*{4}{>{\arraybackslash}X}}
\item $2 \times 5 \times d$ &
\item $3 \times c \times 8$ &
\item $g \times 8 \times 7$ &
\item $3 \times (n + m)$ \\
\item $(a + b) \times 5$ &
\item $b \times ( 5 \times c + 7)$ &
\item $2,5 \times d \times ( d \times 9+ 7 \times 3 )$ &\\
\end{tabularx}
\end{enumerate}
\end{exer1}

\vspace{-1mm}
\begin{exer1}
Donne l'écriture la plus simple de chacun des produits suivants.

\vspace{-2mm}
\begin{enumerate}[font=\color{exer1}]
\begin{tabularx}{\linewidth}{*{6}{>{\arraybackslash}X}}
\item $a \times 1$ &
\item $d \times 0$ &
\item $g \times 1$ &
\item $0 \times c$ &
\item $1 \times b$ &
\item $m \times 1$ \\
\end{tabularx}
\end{enumerate}
\end{exer1}

\vspace{-1mm}
\begin{exer2}
Simplifie les expressions suivantes

\vspace{-2mm}
\begin{enumerate}[font=\color{exer2}]
\begin{tabularx}{\linewidth}{*{5}{>{\arraybackslash}X}}
\item $2 \times a + 5 \times c$ &
\item $a \times d + 5 \times 8$ &
\item $38 \times (3 + 2 \times c)$ &
\item $3 \times z - 0 \times b$ &
\item $3 \times 7 - d \times b$ \\
\item $a \times (3 \times 9 + b \times n)$ &
\item $0 \times u + 1 \times m$ &
\item $a \times 6 \times n + 3 \times p$ &
\item $9 \times m \times 5 + k \times j \times 8$ &\\
\end{tabularx}
\end{enumerate}
\end{exer2}

\vspace{-1mm}
\begin{exer1}
Écris, sans les calculer et en utilisant la notation \og carré \fg{} ou \og cube \fg{}, les produits suivants.

\vspace{-2mm}
\begin{enumerate}[font=\color{exer1}]
\begin{tabularx}{\linewidth}{*{5}{>{\arraybackslash}X}}
\item $6 \times 6$&
\item $n \times n$ &
\item $4 \times 4 \times 4$ &
\item $r \times r \times r$ &
\item $2 \times 2 \times p$ \\
\item $r \times r \times p \times p \times p$&
\item $3 \times 3 \times n \times n$ &
\item $1 \times 1 \times 1 \times y \times y$ &
\item $d \times d \times d \times 6 \times 6$&\\
\end{tabularx}
\end{enumerate}
\end{exer1}

%\subsection*{En fonction de ...}
%\addcontentsline{toc}{subsection}{En fonction de ...}

\vspace{-1mm}
\begin{exer1}
$n$ est un nombre entier. Exprimer en fonction de $n$ :

\vspace{-2mm}
\begin{enumerate}[font=\color{exer1}]
\begin{tabularx}{\linewidth}{*{4}{>{\arraybackslash}X}}
\item la moitié de $n$ &
\item le double de $n$ &
\item le tiers de $n$ &
\item le triple de $n$ \\
\end{tabularx}
\begin{tabularx}{\linewidth}{*{2}{>{\arraybackslash}X}}
\item le nombre entier suivant $n$ &
\item le nombre entier précédent $n$\\
\end{tabularx}
\end{enumerate}
\end{exer1}

\vspace{-1mm}
\begin{exer1}
Si $x$ représente un nombre, comment écrire les expressions suivantes ?

\vspace{-2mm}
\begin{enumerate}[font=\color{exer1}]
\begin{tabularx}{\linewidth}{*{3}{>{\arraybackslash}X}}
\item La somme de $x$ et de $5$ &
\item La différence de $x$ et de $5$ &
\item La différence de 5 et de $x$ \\
\item Le produit de $5$ et de $x$ &
\item La quotient de $x$ et de $5$ & \\
\end{tabularx}
\end{enumerate}
\end{exer1}

\vspace{-1mm}
\begin{exer1}
Relie chaque phrase de gauche à l'expression littérale correspondante de droite.

\renewcommand{\arraystretch}{1.1}
\begin{tabularx}{14cm}{r c X c l}
Somme de $y$ et de $7$ & $\bullet$ & & $\bullet$ & $7 \times (y-3)$ \\
Produit de $7$ par la somme de $y$ et de $3$ & $\bullet$ & & $\bullet$ & $7-y$ \\
Produit de $7$ par la différence entre $y$ et $3$ & $\bullet$ & & $\bullet$ & $y+7\times 3$ \\
Différence du produit de $7$ par $y$ et de $3$ & $\bullet$ & & $\bullet$ & $y+7$ \\
Différence entre $7$ et $y$ & $\bullet$ & & $\bullet$ & $7 \times y + 3$ \\
Somme de $y$ et du produit de $3$ par $7$ & $\bullet$& & $\bullet$ & $7 \times (y+3)$ \\
Somme du produit de $7$ par $y$ et de $3$ & $\bullet$& & $\bullet$ & $7 \times y - 3$ \\
\end{tabularx}
\end{exer1}

\vspace{-1mm}
\begin{exer1}[Écrire en fonction de]

\begin{enumerate}[font=\color{exer1}]
\item Dans une classe de $26$ élèves, on note $x$ le nombre de filles. Exprime le nombre de garçons en fonction de $x$.
\item Sur un parking, il y a $x$ scooters et $y$ voitures. Exprime le nombre de roues en fonctions de $x$ et $y$.
\item On note $c$ le coté d'un carré. Exprime son aire et son périmètre ne fonction de $c$.
\end{enumerate}
\end{exer1}

\vspace{-1mm}
\begin{exer1}
Léa a $18$ billes rouges de plus que de billes noires. On désigne par $x$ le nombre de billes noires.

\begin{enumerate}[font=\color{exer1}]
\item Exprime le nombre de billes rouges en fonctions de $x$.
\item Exprime alors le nombre total de billes en fonction de $x$.
\end{enumerate}
\end{exer1}

\vspace{-1mm}
\begin{exer1}
Dans une assemblée, il y a deux fois plus de Belges que de Luxembourgeois et $48$ Néerlandais de plus que de Luxembourgeois. On désigne par $x$ le nombre de Luxembourgeois. Quelle est la composition de l'assemblée (en fonction de $x$).
\end{exer1}

\vspace{-1mm}
\begin{exer2}
Paul calcule que, s'il achète deux croissants et une brioche à $1,83$ \euro , il dépense $0,47$ \euro ~~ de plus que s'il achète quatre croissants. On désigne par $x$ le prix d'un croissant.

\begin{enumerate}[font=\color{exer2}]
\item Écris, en fonction de $x$, le prix en euros de deux croissants et d'une brioche.
\item Écris le prix en euros de quatre croissants.
\item Écris une égalité.
\end{enumerate}
\end{exer2}

\vspace{-1mm}
\begin{exer1}
Voici trois segments $[AB]$, $[CD]$ et $[EF]$ dont on cherche à calculer les longueurs respectives $AB$, $CD$ et $EF$. Dans chacun des cas, écris une expression permettant de calculer ces longueurs.

\noindent \includegraphics[scale=0.175]{../40.Image/8. Exo-Puissance_de_dix 1}
\end{exer1}

\vspace{-1mm}
\begin{exer1}
Le mur d'une chambre a une forme rectangulaire. Sa longueur mesure le double de sa largeur.

\begin{enumerate}[font=\color{exer1}]
\item Si la largeur mesure $5$m, calcule :

\vspace{-2mm}
\begin{enumerate}
\begin{tabularx}{\linewidth}{*{3}{>{\arraybackslash}X}}
\item[•]la mesure de la longueur &
\item[•]le périmètre du mur &
\item[•]L'aire du mur \\
\end{tabularx}
\end{enumerate}
\item En appelant $\ell$ la mesure en mètres de la largeur, exprime :

\vspace{-2mm}
\begin{enumerate}
\begin{tabularx}{\linewidth}{*{3}{>{\arraybackslash}X}}
\item[•]la mesure de la longueur &
\item[•]le périmètre du mur &
\item[•]L'aire du mur \\
\end{tabularx}
\end{enumerate}
\item Manon refait la tapisserie du mur de sa chambre avec une frise le long du haut du mur. Un rouleau de tapisserie couvre $2$ m$^2$ et un rouleau de frise mesure $50$cm. Exprime en fonction de $\ell$ le nombre de rouleaux de tapisserie et de frise nécessaire.
\end{enumerate}
\end{exer1}

\vspace{-1mm}
\begin{exer1}
La dernière étape du Tour de France $2020$ se composait d'une partie sur la route d'une longueur de $67,5$km et d'un nombre de tours de l'avenue des Champs-Élysée d'une longueur de $7$km.

\begin{enumerate}[font=\color{exer1}]
\item Un coureur cycliste chute au début du $8^e$ tour. Quelle distance a-t-il parcourue ?
\item Soit $t$ le nombre de tours parcourus. Exprime la longueur de l'étape en fonction de $t$.
\item La longueur de l'étape est égale à $122$km. Écris une égalité.
\end{enumerate}
\end{exer1}

\vspace{-1mm}
\begin{exer1}[Avec un programme]

\vspace{5mm}
\noindent \begin{minipage}{5.5cm}
\arrayrulecolor{black}
\renewcommand{\arraystretch}{0.1}
\begin{tabular}{|p{5cm}|}
\hline
\vspace{2mm} 1. \edupythoninline{Choisis un nombre} \\ 
\vspace{2mm} 2. \edupythoninline{Ajoute-lui 2} \\ 
\vspace{2mm} 3. \edupythoninline{Multiplie le resultat par 3} \\ 
\vspace{2mm} \\
\hline
\end{tabular}
\end{minipage}
\hspace{5mm}
\begin{minipage}{10.6cm}
\begin{enumerate}
\item Effectue le programme de calcul pour $5$
\item Soit $x$ le nombre choisi à l'étape $1$. Exprime le résultat du programme de calcul en fonction de $x$.
\end{enumerate}
\end{minipage}
\end{exer1}

\vspace{-1mm}
\begin{exer2}[Avec un autre programme]

\vspace{5mm}
\noindent \begin{minipage}{5.5cm}
\arrayrulecolor{black}
\renewcommand{\arraystretch}{0.1}
\begin{tabular}{|p{5cm}|}
\hline
\vspace{2mm} 1. \edupythoninline{Choisis un nombre} \\ 
\vspace{2mm} 2. \edupythoninline{Calcule son double} \\ 
\vspace{2mm} 3. \edupythoninline{Soustrais 6} \\ 
\vspace{2mm} 4. \edupythoninline{Divise le resultat par 2} \\ 
\vspace{2mm} \\
\hline
\end{tabular}
\end{minipage}
\hspace{5mm}
\begin{minipage}{10.7cm}
\begin{enumerate}
\item Effectue le programme de calcul pour $4$.
\item Recommence avec deux autres nombres. Que constates-tu ?
\item Soit $x$ le nombre choisi à l'étape $1$. Exprime le résultat du programme de calcul en fonction de $x$.
\item Théo choisit un nombre et lui soustrait $3$. Compare les résultats de ces deux programmes.
\end{enumerate}
\end{minipage}
\end{exer2}

%\subsection*{Calculer la valeur d'une expression littérale}
%\addcontentsline{toc}{subsection}{Calculer la valeur d'une expression littérale}

\vspace{-1mm}
\noindent \begin{minipage}{8.5cm}
\begin{exer1}
Calcule la valeur de $B$ et de $Z$ pour $x=5$

\vspace{-2mm}
\begin{enumerate}[font=\color{exer1}]
\begin{tabularx}{\linewidth}{*{2}{>{\arraybackslash}X}}
\item $B = 20x$&
\item $Z = 9x$ \\
\end{tabularx}
\end{enumerate}
\end{exer1}
\end{minipage}
\begin{minipage}{8.7cm}
\begin{exer1}
Calcule la valeur de $M$ et de $A$ pour $y=10$

\vspace{-2mm}
\begin{enumerate}[font=\color{exer1}]
\begin{tabularx}{\linewidth}{*{2}{>{\arraybackslash}X}}
\item $M = 5y+3$&
\item $A = 8y-25$ \\
\end{tabularx}
\end{enumerate}
\end{exer1}
\end{minipage}

\vspace{-1mm}
\begin{exer1}
Calcule la valeur de $T$ et de $Y$ pour $a=10$ et $b=3$

\vspace{-2mm}
\begin{enumerate}[font=\color{exer1}]
\begin{tabularx}{\linewidth}{*{2}{>{\arraybackslash}X}}
\item $T = 7a + 3b - 3$&
\item $Y = 3a - 7b + 4$ \\
\end{tabularx}
\end{enumerate}
\end{exer1}

\vspace{-1mm}
\begin{exer1}
Calcule la valeur de $M$, de $R$ et de $E$ 

\begin{enumerate}[font=\color{exer1}]
\noindent \begin{tabularx}{8cm}{*{2}{>{\arraybackslash}X}}
pour $m=5$ et $n=9$ &
\item $M = 7m + 10n + nm$\\
\item $R = 10n + 5mn -8n$&
\item $E = 8n - 4m - 6mn$ \\
\end{tabularx}
\end{enumerate}
\end{exer1}

\vspace{-1mm}
\begin{exer2}
Complète le tableau suivant :

\vspace{-30mm}
\begin{flushright}
\renewcommand{\arraystretch}{1.7}
\begin{tabularx}{7.5cm}{|>{\columncolor{exer2!20}\centering}p{2.5cm}|*{3}{>{\centering \arraybackslash}X|}}
\hline
\rowcolor{exer2!20} & $x=4$ & $x=0$ & $x=-2$ \\\hline
$3(2x-7)-5x$ & & & \\\hline
$(x-4)(x-2)$ & & & \\\hline
$(2-x)^2$ & & & \\\hline
\end{tabularx}	
\end{flushright}		
\end{exer2}
\end{document}